\section{Policy Interpretation and Research Requirements}\label{sec:policy}

The results motivate scrutiny of the information environment alongside
financial fragility, subject to the predictive validation in
Table~\ref{tab:forecast}. They do not establish that collecting the proposed
index improves decisions, that a financial rumor persists for the aggregate
index's half-life, or that a measured country is near a crisis threshold.

Rapid official communication, message authentication and restrictions on
content propagation are distinct candidate interventions. Their effects may
span the model intercept, persistence, innovation variance, public trust and
the reach of accurate information. Ranking them requires estimates of these
responses, implementation costs and possible adverse effects. The conditional
elasticity result in Proposition~\ref{prop:burden} does not provide those
estimates. Equation~\eqref{eq:cost} makes the missing comparison explicit.

A stronger policy study would link independently identified financial
falsehoods to audience exposure, bank funding or lending outcomes and
credible correction timings. It would then estimate the response to an
intervention using a design with a defensible counterfactual, rather than
relabeling a country-level association. A financial-content validation sample
would also test whether the broad DSP proxy captures the mechanism proposed
here. Out-of-time and real-time validation against richer financial-condition,
uncertainty and geopolitical-risk benchmarks would be needed before proposing
an operational monitoring instrument. These are requirements for stronger
claims, not results claimed by the present paper.
